\section{Conclusion}
\label{sec:conclusion}

% Final numbers land after the 2026-07-19 final; the two remaining
% placeholders are the only edits this section still needs.

We set out to answer a practical question: how much of a fantasy
football tournament can be handled by models that are built, validated,
and revised while the tournament is running? The system that emerged --
three complementary point-prediction backends routed per position by
held-out error, a mixed-integer squad and transfer optimizer that now
prices captaincy ceiling, availability, and the backup-goalkeeper
asymmetry, and a live decision layer audited round by round against its
own recommendations -- answered most of it. Across the seven rounds
completed before the final, every backend beat a constraint-valid
random baseline by a wide margin (\TotalGbm{} points for the strongest
backend and \TotalEnsemble{} for the production ensemble against
\TotalRandomBaselineMean{} for chance), and the human-plus-model entry
scored \TotalUserActualNet{} net points while paying real transfer
costs the replays do not. \POSTFINAL{One sentence completing the record
with the final round: closing cumulative totals and the entry's final
global rank.}

The negative results are, we believe, the most transferable part of
the record. A six-component upgrade lost to the baseline it meant to
replace; a theoretically derived goalkeeper multiplier measured worse
than an empirically calibrated one at less than half its size;
prediction-market prices failed to improve player-level forecasts as a
regression feature, earned only a capped minority weight in match-level
forecasting, and lost to a plain Elo rating out of sample; and minutes
information, whether proxied or real, helped only in the selection
layer, never in the point model. The positive findings mirror them:
leak-free trailing form, tournament retraining, real expected-goals
form, per-position routing, and ceiling-aware captaincy each cleared a
pre-registered walk-forward gate before deployment. None of these
lessons is specific to this World Cup; they are lessons about small
samples, about anchoring on sentiment, and about the discipline of
validating every change on data the model has never seen.

\POSTFINAL{Closing paragraph after the final: the realized outcome of
the locked final-round plan, an honest assessment of the human
overrides against the model recommendations across the tournament, and
the open-source release of the complete system, data, and decision
record.}
